% !TeX program = pdflatex
\documentclass[11pt]{article}

\usepackage[utf8]{inputenc}
\usepackage[T1]{fontenc}
\usepackage[brazilian]{babel}
\usepackage{amsmath,amssymb,amsthm}
\usepackage[margin=1in]{geometry}
\usepackage{booktabs}
\usepackage{hyperref}
\usepackage{microtype}

\hypersetup{colorlinks=true, linkcolor=blue, citecolor=blue, urlcolor=blue}

\newtheorem{theorem}{Teorema}[section]
\newtheorem{lemma}[theorem]{Lema}
\newtheorem{proposition}[theorem]{Proposição}
\newtheorem{corollary}[theorem]{Corolário}
\theoremstyle{definition}
\newtheorem{definition}[theorem]{Definição}
\newtheorem{example}[theorem]{Exemplo}
\newtheorem{problem}[theorem]{Problema}
\theoremstyle{remark}
\newtheorem{remark}[theorem]{Observação}

\newcommand{\Cyc}{\mathrm{Cyc}}
\newcommand{\Per}{\mathrm{Per}}
\newcommand{\dens}{\mathrm{dens}}
\DeclareMathOperator{\ord}{ord}
\DeclareMathOperator{\rad}{rad}
\DeclareMathOperator{\mmc}{mmc}

\title{\textbf{Mapas Iterados de Soma de Dígitos de Potências:\\ um Teorema de Cota Inferior, uma Lei Exata de Densidade de Bacias\\ e a Oscilação da Divisão Intra-Assinatura}}
\author{Alex Martins\\ \small Pesquisa Independente}
\date{Agosto de 2026 (v2, corrigida; ver Apêndice~B)}

\begin{document}
\maketitle

\begin{center}
\small \textbf{MSC 2020:} 37P35, 11A63, 11K06, 11A25, 05C20\\
\small \textbf{Palavras-chave:} sistemas dinâmicos discretos, funções soma de dígitos, números felizes generalizados, aritmética modular, bacias de atração, densidade natural, equidistribuição, grafos funcionais, funções digitais.\\
\small \textbf{Zenodo:} \url{https://zenodo.org/records/22181953} \quad \textbf{DOI:} \url{https://doi.org/10.5281/zenodo.22181953}
\end{center}

\begin{abstract}
Estudo a família de sistemas dinâmicos discretos $f_{k,b}(n) = S_b(n^k)$, em que $S_b$ é a função soma dos dígitos em base $b$ e $k \ge 1$ é um expoente fixo. O mapa relaciona-se com, mas distingue-se das, funções felizes generalizadas $S_{e,b}(n) = \sum_i d_i^{\,e}$ de Grundman e Teeple; a distinção, e a estrutura modular que dela se segue, é precisada na \S1. Toda órbita é eventualmente periódica, e o sistema organiza-se por um esqueleto algébrico: o grafo funcional do mapa de potência modular $\varphi_{k,b-1}(x) = x^k \bmod (b-1)$.

Demonstro três resultados estruturais. \textbf{(I) Teorema da Cota Inferior.} O número de atratores satisfaz $|C(k,b)| \ge \Cyc(\varphi_{k,b-1})$, o número de ciclos do mapa modular. \textbf{(II) Lei de Densidade de Bacias.} Para cada ciclo modular $\gamma_i$ com bacia de resíduos $R_i$, a bacia agregada $\mathcal{N}_i$ de todos os atratores que carregam a assinatura de resíduos $\gamma_i$ possui densidade natural, igual exatamente ao peso modular $p_i = |R_i|/(b-1)$; sobre uma janela $[1,M]$ a proporção empírica desvia no máximo $\min(|R_i|,\,b-2)/M$, e a afirmação subjacente em janela finita é uma identidade exata entre inteiros. \textbf{(III)} Ambas as funções de contagem de $\varphi_{k,m}$ admitem formas fechadas: a contagem de pontos periódicos é a clássica $\prod_{p^e \| m}\big(1 + \kappa_k(\varphi(p^e))\big)$, com $\kappa_k(N)$ o maior divisor de $N$ coprimo com $k$, enquanto a contagem de ciclos (a quantidade de que a cota inferior precisa) obtém-se combinando tipos de ciclo locais pelo teorema chinês dos restos. A distinção importa e passou batida em um rascunho anterior: a contagem de pontos periódicos depende de $k$ apenas via $\rad(k)$, ao passo que a contagem de ciclos não.

Os resultados (I) e (II) são consequências elementares da congruência da soma de dígitos (a ``prova dos noves'') e da equidistribuição de resíduos. A novidade reivindicada restringe-se à organização dessas consequências e ao isolamento do problema restante: como a massa $p_i$ se divide entre vários atratores físicos que compartilham uma assinatura. Essa partição não é elementar. Resolvo-a empírica e mecanicamente. As densidades de bacia individuais não parecem existir; ao longo de comprimento de dígitos fixo elas oscilam quase-periodicamente, em antifase entre atratores competidores, sempre somando $p_i$. Um modelo sem parâmetros, a lei gaussiana de $S_b(n^k)$ (teorema do limite local) convoluída com a rotulação exata inteiro-para-atrator, reproduz a oscilação dentro do ruído de Monte Carlo (MAE $= 0{,}0030$). A oscilação é uma instância da clássica flutuação log-periódica de funções digitais (Delange; Drmota--Grabner).

Todas as afirmações são verificadas computacionalmente: uma varredura exaustiva de $19{,}500$ pares de parâmetros $(k,b)$ com $k \le 500$, $b \le 40$, com zero violações da cota inferior e com a forma inteira exata da lei agregada valendo sem exceção, acrescida de uma reimplementação independente relatada no Apêndice~A.

\smallskip
\noindent\emph{Errata.} Esta versão corrige quatro erros do rascunho de julho de 2026: uma afirmação falsa sobre a contagem de ciclos modulares, uma classificação errada de quais assinaturas oscilam, uma alegação exagerada de estreiteza para a cota de erro, e uma estatística confundida. Cada um está documentado no Apêndice~B, junto com o raciocínio que o produziu, pois os modos de falha são instrutivos.
\end{abstract}

\section{Introdução}

\subsection{O mapa, e como ele difere das funções felizes}

Fixe uma base $b \ge 2$ e ponha $m = b-1$. A \textbf{função feliz generalizada} estudada por Grundman e Teeple \cite{GT01,GT07} e revista em \cite{survey} é
\[
S_{e,b}(n) \;=\; \sum_i d_i^{\,e}, \qquad n = \sum_i d_i\, b^i,
\]
a soma das $e$-ésimas potências dos \emph{dígitos} de $n$ em base $b$. O objeto do presente trabalho é o mapa diferente
\[
f_{k,b}(n) \;=\; S_b\big(n^k\big),
\]
a \textbf{soma dos dígitos de $n^k$} em base $b$. Esses são sistemas dinâmicos distintos. Para $(k,b) = (2,10)$, o mapa feliz envia $13 \mapsto 1^2+3^2 = 10 \mapsto 1$ (logo $13$ é feliz), ao passo que $f_{2,10}(13) = S_{10}(169) = 16$ e $f_{2,10}(16) = S_{10}(256) = 13$: um $2$-ciclo. A iteração clássica dos números felizes e $f_{2,10}$ não são o mesmo sistema.

O que as duas famílias compartilham é a aritmética que dirige tudo o que segue: a congruência clássica
\begin{equation}\label{eq:co9}
S_b(N) \equiv N \pmod{m}, \qquad m = b-1,
\end{equation}
(a ``prova dos noves''), válida para todo $N$, pois $b \equiv 1 \pmod{m}$ implica $d_i b^i \equiv d_i \pmod m$ dígito a dígito. Para $f_{k,b}$ isso produz a identidade
\begin{equation}\label{eq:inv}
f_{k,b}(n) \;=\; S_b(n^k) \;\equiv\; n^k \pmod{m},
\end{equation}
de modo que reduzir a dinâmica mod $m$ fornece exatamente o mapa de potência modular $\varphi_{k,m}(x) = x^k \bmod m$. (Para o mapa feliz $S_{e,b}$ a redução mod $m$ é $\sum_i d_i^{\,e}$, que não é função apenas de $n \bmod m$; é por isso que $f_{k,b}$ possui um esqueleto modular mais simples que a família feliz.) A invariância \eqref{eq:inv} é elementar e clássica. Não se reivindica novidade para ela, apenas para as consequências estruturais e distribucionais reunidas abaixo, e por isolar o problema da divisão da \S\ref{sec:split}.

\subsection{As perguntas}

Um tratamento sistemático da família $f_{k,b}$ parametrizada pelo expoente $k$ e pela base $b$ parece estar ausente da literatura. As perguntas tratadas aqui são:
\begin{enumerate}
\item \textbf{Contagem.} Quantos atratores distintos $f_{k,b}$ possui, e a contagem pode ser limitada a partir apenas dos parâmetros algébricos? (Resposta: o Teorema da Cota Inferior, \S\ref{sec:lb}.)
\item \textbf{Distribuição.} Quanta massa, no sentido de densidade natural, a bacia de cada atrator carrega? (Resposta, no nível agregado: a Lei de Densidade de Bacias, \S\ref{sec:bdl}.)
\item \textbf{Divisão.} Quando vários atratores compartilham uma assinatura de resíduos, como a massa agregada (exatamente determinada) se divide entre eles? (Resposta, empírica e mecanicamente: a divisão não se estabiliza; oscila; \S\ref{sec:split}.)
\end{enumerate}

A instância clássica $(k,b) = (2,10)$ já exibe o fenômeno distribucional. Os três atratores $\{1\}$, $\{9\}$, $\{13,16\}$ capturam os inteiros positivos nas proporções $22{,}2\%,\,33{,}3\%,\,44{,}4\%$, precisamente $2/9,\,3/9,\,4/9$, os tamanhos das classes de resíduos módulo $9$ que alimentam cada atrator. A Seção~\ref{sec:bdl} mostra que isso é forçado por \eqref{eq:inv} juntamente com a equidistribuição, portanto uma lei para a massa agregada, não uma coincidência. A Seção~\ref{sec:split} mostra que a liberdade que a lei deixa intocada, a divisão intra-assinatura, é onde reside a dificuldade restante.

A separação seguinte é usada ao longo de todo o texto:
\begin{itemize}
\item \textbf{Determinado pela aritmética modular (exatamente, elementarmente):} uma cota inferior para o número de atratores, e a massa total de bacia por assinatura de resíduos.
\item \textbf{Não determinado pela aritmética modular:} a contagem excedente de atratores $\Delta(k,b)$, e se a massa de cada assinatura se divide em densidades individuais convergentes.
\end{itemize}

\subsection{Escopo e contribuição}

Os teoremas sobre massa agregada e sobre contagens de ciclos são corolários curtos da congruência da soma de dígitos e da equidistribuição de resíduos. Um leitor fluente na prova dos noves achará as provas breves. O artigo registra quatro coisas.

O Teorema~\ref{thm:lb} dá a cota inferior $|C(k,b)| \ge \Cyc(\varphi_{k,b-1})$. O Teorema~\ref{thm:bdl} dá a densidade agregada de bacia de cada assinatura de resíduos exatamente, $p_i = |R_i|/(b-1)$, com a cota em janela finita $\min(|R_i|,\,b-2)/M$ e com uma identidade inteira exata por baixo (Proposição~\ref{prop:partition}). A Proposição~\ref{prop:cyc} dá uma forma fechada para a contagem de ciclos do mapa modular, que é a quantidade de que o Teorema~\ref{thm:lb} precisa e que não é função de $\rad(k)$. A Observação~\ref{rem:result} registra que, quando vários atratores compartilham uma assinatura, as densidades de bacia individuais não parecem existir: ao longo de comprimento de dígitos fixo as massas oscilam quase-periodicamente, em antifase, sempre somando $p_i$.

A lei agregada torna padrões como $22/33/44$ um teorema, e não folclore, e isola o objeto restante, a divisão intra-assinatura. Essa divisão é uma instância da flutuação log-periódica de funções digitais (Delange \cite{delange}, Drmota--Grabner \cite{drmota-grabner}). A fórmula de pontos periódicos da \S\ref{sec:count} é teoria clássica de digrafos de potências (Somer--K\v{r}\'i\v{z}ek \cite{somer-krizek}, Chou--Shparlinski \cite{chou-shparlinski}). A fórmula de contagem de ciclos da mesma seção é elementar; não a encontrei enunciada para este propósito.

Ambos os teoremas exatos foram checados sem exceção numa grade exaustiva de $19{,}500$ pares $(k,b)$ com $k \le 500$ e $b \le 40$ (Apêndice~A). Esta versão corrige quatro erros do rascunho de julho de 2026 (Apêndice~B).

\section{Definições e Notação}

Fixe $b \ge 2$ e ponha $m = b-1$.

\begin{definition}[Soma de dígitos]
Para $n \ge 1$, $S_b(n)$ é a soma dos dígitos de $n$ na representação usual em base $b$.
\end{definition}

\begin{definition}[O mapa]
O \textbf{mapa iterado de soma de dígitos de potências} é $f_{k,b}\colon \mathbb{Z}^+ \to \mathbb{Z}^+$, $f_{k,b}(n) = S_b(n^k)$.
\end{definition}

\begin{definition}[Atratores]
A \textbf{órbita} de $n$ é $(f_{k,b}^{\,t}(n))_{t \ge 0}$. Um \textbf{atrator} é uma órbita periódica minimal: um conjunto finito $A = \{a_1,\dots,a_L\}$ com $f(a_i) = a_{i+1}$ (índices mod $L$), sem subconjunto próprio periódico. $C(k,b)$ denota o conjunto de todos os atratores de $f_{k,b}$.
\end{definition}

\begin{definition}[Mapa de potência modular]
$\varphi_{k,m}(x) = x^k \bmod m$ sobre $\mathbb{Z}/m\mathbb{Z}$; $G(\varphi_{k,m})$ é seu grafo funcional (todo vértice tem grau de saída $1$, logo cada componente tem forma de $\rho$: um ciclo com árvores entrantes). Seus ciclos são $\gamma_1,\dots,\gamma_c$, $c = \Cyc(\varphi_{k,m})$.
\end{definition}

\begin{definition}[Bacia modular e peso]
Para cada ciclo $\gamma_i$,
\[
R_i = \{\, r \in \mathbb{Z}/m\mathbb{Z} : \varphi_{k,m}^{\,t}(r) \in \gamma_i \text{ para algum } t \ge 0 \,\}.
\]
Por construção $\{0,\dots,m-1\} = \bigsqcup_{i=1}^{c} R_i$. O \textbf{peso modular} de $\gamma_i$ é $p_i = |R_i|/m = |R_i|/(b-1)$; note que $\sum_i p_i = 1$.
\end{definition}

\begin{definition}[Assinatura de resíduos]
Para um atrator físico $A = \{a_1,\dots,a_L\} \in C(k,b)$, sua \textbf{assinatura de resíduos} é
\[
\sigma(A) = \{\, a_j \bmod m : 1 \le j \le L \,\} \subseteq \mathbb{Z}/m\mathbb{Z}.
\]
\end{definition}

\begin{definition}[Bacia e densidade]
$B(A) = \{\, n \in \mathbb{Z}^+ : f_{k,b}^{\,t}(n) \in A \text{ para algum } t \ge 0 \,\}$; sua \textbf{densidade natural}, \emph{quando o limite existe}, é $\delta(A) = \lim_{M\to\infty} |B(A) \cap [1,M]|/M$.
\end{definition}

\begin{definition}[Excesso de bifurcação]
$\Delta(k,b) = |C(k,b)| - \Cyc(\varphi_{k,b-1})$.
\end{definition}

\section{O Esqueleto Modular}

\begin{lemma}[Invariância modular]\label{lem:inv}
Para todo $n \in \mathbb{Z}^+$, $f_{k,b}(n) \equiv n^k \pmod{m}$. Logo, a sequência de resíduos $\big(f_{k,b}^{\,t}(n) \bmod m\big)_{t\ge 0}$ é exatamente a $\varphi_{k,m}$-órbita de $n \bmod m$.
\end{lemma}

\begin{proof}
Por \eqref{eq:co9}, $f_{k,b}(n) = S_b(n^k) \equiv n^k \pmod m$, e indução em $t$ fornece a afirmação sobre a órbita.
\end{proof}

\begin{lemma}[Contração]\label{lem:contraction}
Existe um limiar $N^*(k,b)$ tal que $f_{k,b}(n) < n$ para todo $n > N^*$. Consequentemente, toda órbita entra no conjunto finito $[1, N^*]$ e é eventualmente periódica.
\end{lemma}

\begin{proof}
Se $n$ tem $D$ dígitos em base $b$, então $n \ge b^{D-1}$, enquanto $n^k < b^{kD}$ tem no máximo $kD$ dígitos, cada um no máximo $b-1$, logo $S_b(n^k) \le (b-1)\,k\,D$. Como $b^{D-1}$ cresce exponencialmente em $D$ e $(b-1)kD$ linearmente, existe $D_0$ com $(b-1)kD < b^{D-1}$ para todo $D \ge D_0$; tome $N^* = b^{D_0 - 1}$.
\end{proof}

\begin{corollary}[Região de aprisionamento finita]\label{cor:trap}
Todo atrator de $f_{k,b}$ está em $[1, N^*(k,b)]$, e toda bacia $B(A)$ é determinada pelo primeiro iterado que cai nessa janela. (O minerador da \S\ref{sec:verification} computa a menor cota de ponto fixo desse tipo, $M(k,b)$, e trabalha inteiramente dentro de $[1,M]$.)
\end{corollary}

\section{O Teorema da Cota Inferior}\label{sec:lb}

\begin{theorem}[Teorema da Cota Inferior]\label{thm:lb}
Para todos $k \ge 1$ e $b \ge 2$,
\[
|C(k,b)| \;\ge\; \Cyc\big(\varphi_{k,\,b-1}\big).
\]
\end{theorem}

\begin{proof}
Três passos.

\textbf{Passo 1 (Invariância modular).} Pelo Lema~\ref{lem:inv}, a classe de resíduos de $f_{k,b}(n)$ módulo $m$ é determinada pela de $n$: o mapa $f_{k,b}$ respeita a partição de $\mathbb{Z}^+$ em classes de resíduos mod $m$, e o mapa induzido nos resíduos é exatamente $\varphi_{k,m}$.

\textbf{Passo 2 (Subconjuntos invariantes disjuntos).} Os ciclos de $G(\varphi_{k,m})$ induzem a partição em bacias $\{0,\dots,m-1\} = \bigsqcup_{i=1}^{c} R_i$. Defina
\[
\mathcal{N}_i = \{\, n \in \mathbb{Z}^+ : n \bmod m \in R_i \,\}.
\]
Cada $R_i$ é positivamente invariante sob $\varphi_{k,m}$ (se $r$ flui para $\gamma_i$, o mesmo vale para $\varphi(r)$), logo, pelo Passo~1, cada $\mathcal{N}_i$ é positivamente invariante sob $f_{k,b}$, e os $\mathcal{N}_i$ são dois a dois disjuntos.

\textbf{Passo 3 (Um atrator por subconjunto).} Fixe $i$ e qualquer $n_0 \in \mathcal{N}_i$. Pelo Lema~\ref{lem:contraction}, a órbita de $n_0$ entra no conjunto finito $[1,N^*] \cap \mathcal{N}_i$, que é positivamente invariante pelo Passo~2; sendo finito, a órbita eventualmente se torna periódica. A órbita periódica resultante é um atrator contido em $\mathcal{N}_i$.

\textbf{Conclusão.} Os $\mathcal{N}_1,\dots,\mathcal{N}_c$ são dois a dois disjuntos e cada um contém pelo menos um atrator, logo $|C(k,b)| \ge c = \Cyc(\varphi_{k,b-1})$.
\end{proof}

\begin{remark}
A cota é tipicamente estrita: na verificação da \S\ref{sec:verification}, a igualdade exata $|C| = \Cyc$ é a exceção, não a regra. O excesso $\Delta(k,b) \ge 0$ é um efeito da dinâmica de dígitos (vários atratores em escalas numéricas diferentes compartilhando uma assinatura de resíduos) e é estudado empiricamente nas \S\S\ref{sec:split} e \ref{sec:predictors}.
\end{remark}

\begin{example}[O caso clássico $(k,b) = (2,10)$]\label{ex:classic}
O mapa modular $x \mapsto x^2 \bmod 9$ sobre $\{0,\dots,8\}$:
\[
0\to 0,\quad 1\to 1,\quad 2\to 4,\quad 3\to 0,\quad 4\to 7,\quad 5\to 7,\quad 6\to 0,\quad 7\to 4,\quad 8\to 1,
\]
com ciclos $\{0\},\{1\},\{4,7\}$ e bacias $R = \{0,3,6\},\{1,8\},\{2,4,5,7\}$. Os atratores físicos são $\{1\},\{9\},\{13,16\}$, exatamente $3 = \Cyc(\varphi_{2,9})$, logo $\Delta(2,10) = 0$: este é um dos raros sistemas com igualdade exata.
\end{example}

\begin{example}[Um caso com bifurcação $(k,b) = (3,10)$]\label{ex:bif}
O mapa modular sobre $\mathbb{Z}/9\mathbb{Z}$ tem três pontos fixos $\{0\},\{1\},\{8\}$, prevendo $\ge 3$ atratores. O sistema físico tem sete: $\{1\},\{8\},\{17\},\{18\},\{19,28\},\{26\},\{27\}$. Só dentro da assinatura $\{0\}$:
\[
9 \to S_{10}(729) = 18 \to S_{10}(5832) = 18 \ (\text{fixo}), \qquad 27 \to S_{10}(19683) = 27 \ (\text{fixo}).
\]
Tanto $18$ quanto $27$ são $\equiv 0 \pmod 9$, mas são pontos fixos distintos em escalas diferentes: a teoria modular conta um ciclo para a classe $0$, enquanto a dinâmica real bifurca. Logo $\Delta(3,10) = 7 - 3 = 4$.
\end{example}

\section{A Lei de Densidade de Bacias (Forma Agregada)}\label{sec:bdl}

\begin{lemma}[Determinacy da assinatura]\label{lem:sig}
Para todo $n \in \mathbb{Z}^+$, o atrator $A(n)$ eventualmente atingido pela órbita de $n$ tem assinatura de resíduos
\[
\sigma\big(A(n)\big) = \gamma_i, \quad \text{onde } i \text{ é o único índice com } (n \bmod m) \in R_i.
\]
\end{lemma}

\begin{proof}
Pelo Lema~\ref{lem:inv} a sequência de resíduos da órbita é a $\varphi_{k,m}$-órbita de $r = n \bmod m$, que entra no ciclo $\gamma_i$ com $r \in R_i$. Uma vez que a órbita física está no atrator $A(n)$ (um ciclo de período $L$), seus resíduos são $\varphi$-periódicos; percorrer $A(n)$ uma vez retorna ao início, logo o período de resíduos divide $L$, e como $\varphi$ restrito a um ciclo é uma permutação cíclica desse ciclo, os resíduos traçados são o ciclo inteiro $\gamma_i$. Logo $\sigma(A(n)) = \gamma_i$.
\end{proof}

\begin{proposition}[Partição exata de massa em janela finita]\label{prop:partition}
Fixe $(k,b)$. Para todo $M \ge 1$,
\[
\sum_{\substack{A \in C(k,b) \\ \sigma(A) = \gamma_i}} \big|\, B(A) \cap [1,M] \,\big| \;=\; \big|\, \mathcal{N}_i \cap [1,M] \,\big|.
\]
\end{proposition}

\begin{proof}
As bacias $\{B(A)\}_{A \in C(k,b)}$ particionam $\mathbb{Z}^+$ (toda órbita atinge um único atrator, pelo Lema~\ref{lem:contraction}). Pelo Lema~\ref{lem:sig}, $n \in \mathcal{N}_i$ sse $\sigma(A(n)) = \gamma_i$; logo a união das bacias com assinatura $\gamma_i$ é exatamente $\mathcal{N}_i$. Intersecte com $[1,M]$ e conte.
\end{proof}

\begin{theorem}[Lei de Densidade de Bacias, forma agregada]\label{thm:bdl}
Para todos $k \ge 1$, $b \ge 2$, e todo ciclo modular $\gamma_i$ de $\varphi_{k,b-1}$, a bacia agregada $\mathcal{N}_i = \bigcup_{\sigma(A) = \gamma_i} B(A)$ possui densidade natural, igual ao peso modular exatamente:
\[
q_i \;:=\; \delta(\mathcal{N}_i) \;=\; p_i \;=\; \frac{|R_i|}{b-1}.
\]
Além disso, escrevendo $M = qm + s$ com $0 \le s < m$, a proporção empírica $\hat q_i(M) = |\mathcal{N}_i \cap [1,M]|/M$ satisfaz
\[
\big|\, \hat q_i(M) - p_i \,\big| \;\le\; \frac{\min\big(|R_i|,\; m-|R_i|,\; s\big)}{M} \;\le\; \frac{\min(|R_i|,\, b-2)}{M} \;\le\; \frac{b-1}{M}.
\]
\end{theorem}

\begin{proof}
Pela Proposição~\ref{prop:partition}, $\hat q_i(M) = |\mathcal{N}_i \cap [1,M]|/M$, e $\mathcal{N}_i \cap [1,M]$ é o conjunto dos inteiros em $[1,M]$ cujo resíduo mod $m$ está em $R_i$. Entre $\{1,\dots,M\}$, exatamente $s$ classes de resíduos (chame esse conjunto de $S$) contêm $q+1$ elementos e as restantes $m-s$ contêm $q$. Logo
\[
\big| \mathcal{N}_i \cap [1,M] \big| = |R_i| \cdot \frac{M}{m} + \theta, \qquad \theta = \sigma_i - \frac{s\,|R_i|}{m}, \quad \sigma_i = |R_i \cap S|.
\]
Como $0 \le \sigma_i \le \min(|R_i|, s)$ e $\theta$ é uma combinação convexa limitada por ambos os extremos, $|\theta| \le \min(|R_i|,\, s)$. Aplicando o mesmo argumento ao complementar $\mathbb{Z}/m\mathbb{Z} \setminus R_i$, cujo desvio é $-\theta$, obtém-se também $|\theta| \le m - |R_i|$. Dividindo por $M$ e fazendo $M \to \infty$ resulta $\delta(\mathcal{N}_i) = p_i$.
\end{proof}

\begin{remark}[A cota não é $(b-1)/M$]\label{rem:sharp}
A forma grosseira $|\theta| \le |R_i| \le m$ dá $(b-1)/M$, que é o que um rascunho anterior relatou e descreveu como saturada. Ela não é: como $s \le m-1$ e $\min(|R_i|, m-|R_i|) \le \lfloor m/2 \rfloor$, a constante $b-1$ nunca é atingida, e ao longo da varredura da \S\ref{sec:verification} a pior razão observada entre erro real e $(b-1)/M$ é $0{,}25$. A constante estreita $\min(|R_i|, b-2)$ é aproximada, com razão observada $0{,}95$. A consequência prática é que uma verificação relatando ``$100\%$ das medições dentro de $(b-1)/M$'' testa muito pouco; a afirmação com conteúdo é a identidade inteira da Proposição~\ref{prop:partition}, que não tem folga.
\end{remark}

\begin{remark}[Escopo do Teorema~\ref{thm:bdl}]
O teorema afirma a densidade da união $\mathcal{N}_i$. Ele não afirma que as bacias individuais $B(A)$ dos vários atratores que compartilham a assinatura $\gamma_i$ possuam, cada uma, densidade natural. Quando o excesso local de bifurcação é zero (um atrator por assinatura), as duas coisas coincidem trivialmente. Quando é positivo, a decomposição $q_i = \sum_{\sigma(A)=\gamma_i} \delta(A)$ pressupõe que os limites individuais existam, o que é falso em geral, como mostra a \S\ref{sec:split}: as densidades individuais oscilam e não convergem, enquanto sua soma permanece exatamente $p_i$. (Uma família finita de densidades oscilantes pode somar um conjunto com densidade exata.) Na família vizinha dos números felizes, a densidade natural do conjunto feliz também é conhecida por não convergir (Gilmer \cite{gilmer}).
\end{remark}

\begin{corollary}[Conservação de massa entre espaços]\label{cor:conservation}
A aritmética modular determina a massa total de bacia de cada assinatura de resíduos exatamente, independentemente do excesso de bifurcação $\Delta(k,b)$. O excesso governa apenas a partição interna de cada massa $p_i$ entre os $\ge 1$ atratores físicos de assinatura $\gamma_i$, inclusive se essa partição está bem definida como um vetor de densidades.
\end{corollary}

\begin{example}[A lei dos $22/33/44$]\label{ex:223344}
Para $(k,b) = (2,10)$, o Teorema~\ref{thm:bdl} com as bacias do Exemplo~\ref{ex:classic} prevê densidades agregadas $3/9, 2/9, 4/9$ para as assinaturas $\{0\},\{1\},\{4,7\}$, realizadas pelos atratores $\{9\},\{1\},\{13,16\}$. Medido sobre $[1,3000]$: $33{,}33\% / 22{,}23\% / 44{,}43\%$, dentro da cota $9/3000 = 0{,}003$ dos valores exatos, como o teorema exige.
\end{example}

\section{Pontos periódicos e ciclos do mapa de potência modular}\label{sec:count}

O Teorema~\ref{thm:lb} limita inferiormente por $\Cyc(\varphi_{k,b-1})$, de modo que o lado modular do problema de contagem reduz-se a duas quantidades: quantos pontos são periódicos sob $x \mapsto x^k$ módulo $m$, e em quantos ciclos eles se organizam. Ambas admitem formas fechadas. A distinção entre elas é essencial e foi confundida em um rascunho anterior deste trabalho (Apêndice~B.1).

\begin{proposition}[Contagem de pontos periódicos]\label{prop:per}
Para $k \ge 2$ e $m \ge 2$, o número de pontos periódicos de $x \mapsto x^k$ em $\mathbb{Z}/m\mathbb{Z}$ é
\[
\#\Per(k,m) \;=\; \prod_{p^{e} \,\|\, m} \Big( 1 + \kappa_k\big(\varphi(p^{e})\big) \Big),
\]
onde $\varphi$ é a totiente de Euler e $\kappa_k(N)$ é o maior divisor de $N$ coprimo com $k$.
\end{proposition}

\begin{proof}
Pelo teorema chinês dos restos basta tratar $\mathbb{Z}/p^e\mathbb{Z}$. Um não-unidade não nulo é nilpotente sob $x \mapsto x^k$ para $k \ge 2$: sua valuação $p$-ádica multiplica-se por $k$ a cada passo e atinge $\ge e$, logo flui para $0$; assim $0$ é o único não-unidade periódico. No grupo de unidades $(\mathbb{Z}/p^e\mathbb{Z})^\times$, de ordem $\varphi(p^e)$, um elemento $x$ é periódico sob $g \mapsto g^k$ sse $x^{k^t} = x$ para algum $t \ge 1$, sse $\ord(x) \mid k^t - 1$, sse $\gcd(\ord(x), k) = 1$. O conjunto de tais elementos é o único subgrupo de Hall de $k'$-torção, cuja ordem é o maior divisor de $\varphi(p^e)$ coprimo com $k$, a saber $\kappa_k(\varphi(p^e))$. Acrescentando o ponto $0$ obtém-se o fator local $1 + \kappa_k(\varphi(p^e))$; os fatores multiplicam-se sobre os primos pelo TCR.
\end{proof}

\begin{example}
Para $(k,m) = (2,9)$: $\varphi(9) = 6$, $\kappa_2(6) = 3$, logo $\#\Per = 1 + 3 = 4$, a saber $\{0,1,4,7\}$, conferindo com o Exemplo~\ref{ex:classic}.
\end{example}

Checado contra enumeração por força bruta com 0 discrepâncias em todos os $k \in [2,79]$, $m \in [2,300)$ (\texttt{scripts/cycle\_structure.py}), e reverificado independentemente em $k \in [2,39]$, $m \in [2,199)$ (Apêndice~A). A fórmula é clássica na teoria de digrafos de potências $x \mapsto x^k \bmod n$ (Somer--K\v{r}\'i\v{z}ek \cite{somer-krizek}; Chou--Shparlinski \cite{chou-shparlinski}); inclui-se para fixar notação e por ser o insumo da Proposição~\ref{prop:cyc}, não como resultado novo.

\begin{corollary}[Dependência radical da contagem de pontos periódicos]\label{cor:radical}
$\#\Per(k,m)$ depende de $k$ apenas via $\rad(k)$, o conjunto dos primos que dividem $k$; em particular sua parte $2$-ádica colapsa inteiramente na única transição $2 \nmid k \to 2 \mid k$ e permanece inalterada para todas as valuações $2$-ádicas superiores $v_2(k)$.
\end{corollary}

\begin{proof}
Imediato da Proposição~\ref{prop:per}, pois $\kappa_k(N)$ retira de $N$ exatamente os primos que dividem $k$.
\end{proof}

Essa propriedade não se transfere para a contagem de ciclos. O Corolário~\ref{cor:radical} trata de pontos periódicos; o Teorema~\ref{thm:lb} precisa de $\Cyc$. Os dois não são intercambiáveis, pois um conjunto de pontos periódicos de tamanho conhecido pode organizar-se em números muito diferentes de ciclos. A afirmação correta do lado dos ciclos é a seguinte.

\begin{proposition}[Contagem de ciclos]\label{prop:cyc}
Sejam $k \ge 2$ e $m \ge 2$. Para cada potência prima $p^e \,\|\, m$ seja $H_{p^e} \le (\mathbb{Z}/p^e\mathbb{Z})^\times$ o subgrupo dos elementos cuja ordem é coprima com $k$, e seja $c_d$ o número de elementos de $H_{p^e}$ de ordem exatamente $d$. O tipo de ciclo local de $\varphi_{k,p^e}$ sobre seu conjunto periódico é o multiconjunto
\[
T_{p^e} \;=\; \{1\} \;\cup\; \bigcup_{d \,\mid\, \exp H_{p^e}} \Big\{\, \underbrace{\ord_d(k), \dots, \ord_d(k)}_{c_d/\ord_d(k) \ \text{vezes}} \,\Big\},
\]
onde $\ord_d(k)$ é a ordem multiplicativa de $k$ módulo $d$ (e $\ord_1(k) = 1$), e o singleton $\{1\}$ é o ponto fixo $0$. Então
\[
\Cyc(\varphi_{k,m}) \;=\; \sum_{(\ell_1,\dots,\ell_t) \,\in\, T_{p_1^{e_1}} \times \cdots \times T_{p_t^{e_t}}} \frac{\ell_1 \cdots \ell_t}{\mmc(\ell_1,\dots,\ell_t)}.
\]
\end{proposition}

\begin{proof}
Pela Proposição~\ref{prop:per} o conjunto periódico de $\varphi_{k,p^e}$ é $\{0\} \sqcup H_{p^e}$, e $\varphi_{k,m}$ restrito a ele é uma bijeção. Um elemento $x \in H_{p^e}$ de ordem $d$ satisfaz $\varphi^{\,t}(x) = x^{k^t} = x$ sse $k^t \equiv 1 \pmod d$, logo seu comprimento de ciclo é exatamente $\ord_d(k)$; os $c_d$ elementos de ordem $d$ preenchem portanto $c_d/\ord_d(k)$ ciclos. Pelo teorema chinês dos restos o conjunto periódico de $\varphi_{k,m}$ é o produto dos locais e $\varphi$ age coordenada a coordenada. Uma $t$-upla de ciclos locais de comprimentos $\ell_1,\dots,\ell_t$ abrange $\prod \ell_i$ pontos periódicos, sobre os quais a permutação produto tem ordem $\mmc(\ell_i)$ e age com todas as órbitas desse comprimento comum; logo decompõe-se em $\prod \ell_i / \mmc(\ell_i)$ ciclos. Somando sobre as $t$-uplas obtém-se a fórmula.
\end{proof}

Checado contra enumeração por força bruta com 0 discrepâncias para todos os $k \in [1,60]$, $m \in [1,259]$ (\texttt{scripts/cycle\_structure.py}, \texttt{tests/test\_modular.py}).

\begin{corollary}[A contagem de ciclos não é função de $\rad(k)$]\label{cor:notradical}
$\Cyc(\varphi_{k,m})$ depende de $k$ através das ordens multiplicativas $\ord_d(k)$, logo de $k$ módulo as ordens dos elementos $d$, e não apenas dos primos que dividem $k$.
\end{corollary}

A dependência não é um efeito marginal. Para $m = 41$ o conjunto periódico é $\{0\}$ juntamente com o subgrupo de ordem $5$, sobre o qual $x \mapsto x^k$ age como multiplicação por $k$ módulo $5$; como $\ord_5(2) = 4$, $\ord_5(4) = 2$ e $16 \equiv 1 \pmod 5$,
\[
\Cyc(\varphi_{2,41}) = 3, \qquad \Cyc(\varphi_{4,41}) = 4, \qquad \Cyc(\varphi_{16,41}) = 6,
\]
enquanto $\#\Per = 6$ em todos os casos. Analogamente $\Cyc(\varphi_{2,37}) = 4$ mas $\Cyc(\varphi_{4,37}) = 6$, com $\#\Per = 10$ nos dois casos. Em $m \in [2,500)$ há $421$ módulos com $\Cyc(\varphi_{2,m}) \ne \Cyc(\varphi_{4,m})$, todos eles com $\#\Per$ idêntico, e $21$ dos $38$ módulos $m = b-1 \le 39$ dentro da grade de varredura da \S\ref{sec:verification} se comportam assim.

\begin{remark}[Escopo]
A Proposição~\ref{prop:cyc} fecha completamente o lado modular da contagem: tanto $\#\Per$ quanto $\Cyc$ são agora computáveis sem construir grafo algum. Ela nada diz sobre o excesso físico $\Delta = |C| - \Cyc$, cujos atratores extras são um efeito da dinâmica de dígitos e não modular; isso permanece aberto (Problema~\ref{prob:excess}). Tampouco autoriza qualquer afirmação de que $\Cyc$ seja monótona ou plana em $v_2(k)$; ver Apêndice~B.1 para o enunciado retratado.
\end{remark}

\section{Oscilação da divisão intra-assinatura}\label{sec:split}

O Teorema~\ref{thm:bdl} fixa cada massa agregada $p_i$ exatamente. A pergunta restante é como $p_i$ se divide entre os atratores físicos que compartilham $\gamma_i$.

\subsection{Formulação}

Fixe $(k,b)$ e uma assinatura $\gamma_i$ que hospeda atratores físicos $A_1,\dots,A_r$ ($r = 1 + {}$excesso local), cujas bacias particionam $\mathcal{N}_i$. Duas perguntas:
\begin{itemize}
\item \textbf{(Existência)} Cada $\delta(A_j) = \lim_M |B(A_j) \cap [1,M]|/M$ existe?
\item \textbf{(Valor)} Se sim, qual é o vetor $(\delta(A_1),\dots,\delta(A_r))$? Pela Proposição~\ref{prop:partition} ele deve satisfazer $\sum_j \delta(A_j) = p_i$.
\end{itemize}

\subsection{Redução a uma distribuição de soma de dígitos}

Pelo Corolário~\ref{cor:trap}, todo atrator está em $[1,M(k,b)]$, e um iterado já pousa a órbita perto dessa escala: $f_{k,b}(n) = S_b(n^k) \le (b-1)(1 + k \log_b n)$. Seja $F = f_{k,b}\big|_{[1,M]}$ o mapa finito restrito; cada $A_j$ tem uma bacia finita $\beta_j = B(A_j) \cap [1,M]$, e os $\beta_j$ particionam $[1,M]$ dentro da assinatura $\gamma_i$. Como a pertinência à bacia de $n$ é decidida por onde cai o primeiro iterado,
\[
B(A_j) = \{\, n : f_{k,b}^{\,t}(n) \in \beta_j \text{ para o primeiro } t \text{ com } f^{\,t}(n) \le M \,\},
\]
logo, rastreando o sítio de aterrissagem de primeira passagem, a densidade individual seria
\[
\delta(A_j) \;=\; \sum_{v \in \beta_j} \dens\big\{\, n : \text{o primeiro iterado de } n \text{ em } [1,M] \text{ é igual a } v \,\big\},
\]
sempre que as densidades do lado direito existam. O problema fica assim reduzido à distribuição assintótica da soma de dígitos $S_b(n^k)$ (e de seus primeiros iterados) quando $n$ percorre uma classe de resíduos, um objeto bem estudado.

\subsection{Controle analítico disponível}

A distribuição de somas de dígitos de sequências polinomiais é um assunto maduro. Para $k=2$, Mauduit e Rivat \cite{mauduit-rivat} provaram que $S_b(n^2)$ é assintoticamente equidistribuída em classes de resíduos e satisfaz um teorema central do limite; Drmota--Mauduit--Rivat \cite{dmr} dão teoremas do limite local para $S_b(n^k)$; Drmota \cite{drmota-poly} trata a função somatória. Esses resultados controlam a medida empurrada para frente acima, e também preveem flutuações, que se revelam o ponto essencial.

\subsection{Resolução empírica: a divisão oscila}

Medição de Monte Carlo da divisão restrita a $n$ com exatamente $D$ dígitos em base $b$ (\texttt{scripts/split\_scale.py}, $(k,b) = (3,10)$, $4 \le D \le 90$, $60{,}000$ amostras por faixa, amostragem $\sigma \approx 0{,}002$) dá uma resposta nítida e reproduzível:
\begin{itemize}
\item A massa agregada por assinatura é $1/3$ em todo comprimento de dígitos $D$ (Teorema~\ref{thm:bdl}, como esperado).
\item Nenhuma assinatura tem divisão convergente; todas as sete curvas são não monótonas. Para a assinatura $\{0\}$ os dois pontos fixos $\{18\}$ e $\{27\}$ trocam massa em antifase, quase-periodicamente em $D$, nas faixas $[0{,}060,0{,}239]$ e $[0{,}091,0{,}269]$, amplitude $\approx 0{,}179$, sempre somando $1/3$. Esta é a instância maior e mais visível.
\item As outras duas assinaturas oscilam também, com amplitude menor mas o mesmo comportamento qualitativo. Dentro da assinatura $\{1\}$ o $2$-ciclo $\{19,28\}$ carrega a maior parte da massa mas varia em $[0{,}266,0{,}338]$; dentro da assinatura $\{8\}$ o ponto fixo $\{26\}$ varia em $[0{,}227,0{,}337]$.
\item Os três pontos fixos pequenos $\{1\},\{8\},\{17\}$ não decaem a zero. Cada um permanece em ou perto de zero por longos trechos de $D$ e depois retorna: $\{17\}$ atinge $0{,}100$ em $D = 10$, é numericamente zero ao longo de $22 \le D \le 70$, e volta a $0{,}027$ em $D = 79$; $\{8\}$ é zero ao longo de $13 \le D \le 49$ e tem pico $0{,}028$ em $D = 58$; $\{1\}$ atinge $0{,}068$ em $D = 7$, some em $13 \le D \le 64$, e retorna a $0{,}026$ em $D = 73$. Os retornos estão duas ordens de magnitude acima do ruído de amostragem $\sigma \approx 0{,}002$.
\end{itemize}

\textbf{Mecanismo.} Na redução acima o sítio de aterrissagem é governado pelos iterados $m_1 = S_b(n^k) \approx \tfrac{b-1}{2} k D$, que cresce linearmente em $D$, e em seguida $m_2 = S_b(m_1^{\,k})$, cuja média cresce apenas como $\log D$. Como a região de aprisionamento $[1,M]$ é uma janela fixa (aqui $M = 57$), a distribuição de aterrissagem de primeira passagem vive num conjunto finito fixo e nunca escapa ao infinito; ela apenas deriva e se espalha lentamente. Ao varrer a rotulação fixa de bacias dos inteiros, bacias comparáveis alternam a captura, e bacias pequenas situadas na parte baixa da janela são revisitadas sempre que a cascata de vários passos produz um valor de aterrissagem baixo. As bacias relevantes dentro de $[1,57]$ são
\[
\beta_{\{1\}} = \{1,4,7,10,40\}, \quad \beta_{\{8\}} = \{2,5,8,11,20,50\}, \quad \beta_{\{17\}} = \{14,17,23,47\},
\]
cada uma pequena mas contendo um elemento bem acima dos demais; esses elementos altos são o que a distribuição de aterrissagem em deriva varre de volta em $D$ grande. É por isso que os pontos fixos pequenos recorrem em vez de desaparecer.

\begin{remark}[Nota metodológica]
Uma versão anterior desta seção relatou que as assinaturas $\{1\}$ e $\{8\}$ convergem. Essa conclusão veio de uma regra de veredito que examinava a deriva da divisão ao longo de um único passo de $D$ e declarava estabilização quando ela caía abaixo do ruído de amostragem. Uma curva quase-periódica tem trechos longos e planos, de modo que tal regra reporta convergência sempre que a janela amostrada acontece de estar dentro de um deles. A medição corrigida relata a curva inteira e a amplitude sobre a faixa completa de $D$, e nunca infere convergência a partir de uma inclinação local. Ver Apêndice~B.2.
\end{remark}

\subsection{Um modelo preditivo sem parâmetros}

A oscilação é explicada por um modelo sem parâmetros (\texttt{scripts/split\_predict.py}): modele $m_1 = S_b(n^k)$ como gaussiana de média $\tfrac{b-1}{2}L$ e variância $L\,\tfrac{b^2-1}{12}$ (o número de dígitos $L$ sendo uma mistura computável sobre a faixa de $D$ dígitos; este é o teorema do limite local de \cite{mauduit-rivat,dmr}), restrinja à classe de resíduos da assinatura, convolua com a rotulação exata inteiro-para-atrator $a(v)$ e escalone por $p_i$. Para a assinatura $\{0\}$ de $(3,10)$ isso reproduz a curva medida $\delta_j(D)$ com MAE $= 0{,}0030$ ao longo de $4 \le D \le 60$, sem nenhum parâmetro ajustado. A oscilação é portanto a varredura gaussiana da soma de dígitos $m_1$ sobre a rotulação fixa de bacias dos inteiros.

Com $40{,}000$ amostras por faixa o piso de ruído da medição é $0{,}5/\sqrt{40000} = 0{,}0025$, de modo que um MAE de $0{,}0030$ diz que o modelo é compatível com os dados. Isso não o destaca frente a outros candidatos que também reproduzam fase e amplitude. O conteúdo do modelo é mecanístico, não estatístico: ele prevê a oscilação a partir do teorema do limite local mais a rotulação exata, sem nada ajustado. Distingui-lo de alternativas exigiria ou uma amostra muito maior ou uma previsão da estrutura de ordem superior, que é o que o Problema~\ref{prob:split} pede.

\begin{remark}[Observação empírica, com um mecanismo]\label{rem:result}
Nenhuma densidade de bacia individual $\delta(A_j)$ parece existir: ao longo de comprimento de dígitos fixo as massas oscilam quase-periodicamente em $D$, para toda assinatura com excesso local positivo. Só o agregado por assinatura tem densidade, igual a $p_i$ (Teorema~\ref{thm:bdl}).

Isto é registrado como observação e não como resultado: a não existência de um limite é uma afirmação sobre todo $D$, e um número finito de faixas não a estabelece. O que os dados estabelecem é que as massas individuais não se assentaram até $D = 90$, com amplitudes duas ordens de magnitude acima do ruído de amostragem. O passo da oscilação em $D$ para a não existência da densidade cumulativa é elementar. Escrevendo $\delta_j(d)$ para a massa de $A_j$ entre os inteiros de $d$ dígitos, a proporção cumulativa até $M = b^D$ é uma média ponderada $\sum_{d \le D} w_d \, \delta_j(d)$ com $w_d \propto b^{d}$, de modo que a faixa $D$ sozinha carrega peso $1 - 1/b$. Uma oscilação persistente de amplitude $a$ em $\delta_j(D)$ sobrevive portanto na média cumulativa com amplitude pelo menos $(1-1/b)\,a$ menos a contribuição da cauda geometricamente amortecida. Para $(3,10)$ com $a \approx 0{,}18$ isso é grande demais para ser lavado.
\end{remark}

\subsection{Relação com a teoria das funções digitais}

Não encontrei essa oscilação da divisão de bacias escrita à letra para o mapa $S_b(n^k)$. Ela não é um tipo novo de fenômeno: é uma manifestação da clássica flutuação log-periódica de funções digitais (Delange \cite{delange}, 1975), e a distribuição dirigente $m_1 = S_b(n^k)$ é exatamente o objeto do trabalho de Drmota sobre somas de dígitos de sequências polinomiais \cite{drmota-poly} e da monografia de Drmota--Grabner \cite{drmota-grabner}. O modelo de varredura gaussiana acima é a maquinaria padrão de Mellin--Perron / somas de dígitos polinomiais dessa teoria, aplicada numericamente a este mapa. Um especialista consideraria a oscilação esperada e as ferramentas para torná-la rigorosa, padrão. O que se reivindica aqui é portanto: (i) uma lei agregada exata, elementar e computacionalmente verificada (Teorema~\ref{thm:bdl}); (ii) uma descrição empírica e modelada de por que as densidades individuais deixam de existir para esta família, realizada como instância de oscilação de funções digitais. O próximo passo é verificar \cite{delange,drmota-poly,drmota-grabner} se o enunciado específico da divisão de bacias já está implícito e, se uma nota se justificar, derivar o termo exato de oscilação de Fourier/fractal via Mellin--Perron em vez da gaussiana numérica.

\section{Verificação Computacional}\label{sec:verification}

\subsection{Método}

Uma varredura exaustiva, finita por construção (\texttt{scripts/sweep.py}). Para cada $(k,b)$ ele computa a cota de contração $M(k,b)$, o ponto fixo de $S_b(n^k) \le (b-1)\cdot\mathrm{digits}_b(n^k)$, que pelo Lema~\ref{lem:contraction} garante que todo atrator está em $[1,M]$. Constrói o grafo funcional completo de $f_{k,b}$ sobre $[1,M]$, extrai todos os atratores, bacias e assinaturas de resíduos $\sigma(A) \bmod (b-1)$; um motor independente reconstrói $G(\varphi_{k,b-1})$ e os pesos $p_i$. Precisão arbitrária via \texttt{gmpy2}; a varredura é paralelizada em 28 processos.

\subsection{Resultados da varredura exaustiva}

\begin{center}
\begin{tabular}{@{}ll@{}}
\toprule
\textbf{Parâmetro} & \textbf{Valor} \\
\midrule
Expoentes $k$ & $1$ a $500$ \\
Bases $b$ & $2$ a $40$ \\
Total de pares varridos (exaustivo) & $19{,}500$ \\
Violações da cota inferior ($|C| < \Cyc$) & $0$ \\
Falhas da identidade inteira exata (Prop.~\ref{prop:partition}) & $0$ \\
Falhas do Lema~\ref{lem:sig} (assinatura $=$ ciclo modular completo) & $0$ \\
Comparações de densidade de assinatura & $152{,}276$ \\
Pontos com erro $\le \min(|R_i|,b-2)/M$ (cota estreita) & $152{,}276$ ($100{,}00\%$) \\
Erro absoluto médio $|\hat q_i - p_i|$ & $0{,}0001$ \\
Pior erro isolado & $0{,}10$ em $(k,b,M) = (1,3,5)$ \\
Máximo excesso observado & $\Delta = 98$ em $(k,b) = (451,32)$ \\
\bottomrule
\end{tabular}
\end{center}

As três primeiras linhas são as substantivas: são afirmações inteiras exatas, de modo que um único contraexemplo em qualquer ponto da grade falsificaria um teorema. As linhas de densidade são mais fracas do que parecem. A cota $(b-1)/M$ do Teorema~\ref{thm:bdl} não é aproximada (ao longo da grade a pior razão entre erro real e cota é $0{,}25$; Observação~\ref{rem:sharp}), de modo que ``$100{,}00\%$ dentro da cota'' é quase automático e não deve ser lido como um teste rigoroso. Relatar $152{,}276$ comparações infla a evidência aparente pela mesma razão: as comparações não são independentes, pois todas as assinaturas de um par compartilham uma janela.

A forma estreita é a identidade inteira da Proposição~\ref{prop:partition}: para todo $(k,b)$ e toda janela $[1,M]$, o número de $n \le M$ cuja órbita atinge um atrator de assinatura $\gamma_i$ é igual, exatamente, ao número de $n \le M$ com $n \bmod (b-1) \in R_i$. Nenhum ponto flutuante está envolvido e não há folga para absorver um erro. Essa identidade valeu sem exceção ao longo da grade, e é ela pelo que a verificação deve ser julgada.

Isto verifica a lei agregada; nada diz sobre a divisão, que a \S\ref{sec:split} resolve à parte.

\subsection{Verificação amostrada anterior (superada)}

Uma varredura protótipo anterior testou $1{,}990$ pares amostrados (uma amostra incremental, não uma grade exaustiva) com $k$ atingindo $158$ e $b$ atingindo $20$, usando uma amostra finita $n \le 500$ por célula: zero violações, com igualdade exata $|C| = \Cyc$ em apenas $2$ de $1{,}990$ pares ($\approx 0{,}1\%$). Uma reamostragem independente mais densa sobre $k \in [1,25]$, $b \in [2,12]$ (275 pares, $N = 300$) encontrou 0 violações e uma taxa de igualdade exata de $7{,}64\%$. A frequência de igualdade exata é sensível à faixa amostrada, ela própria uma questão em aberto (Problema~\ref{prob:tight}). Essas cifras são superadas pela varredura exaustiva acima, mas são relatadas por proveniência.

\subsection{Fenomenologia do mapa de calor de bifurcação}

Uma renderização sistemática de $|C(k,b)|$ sobre a grade $(k,b)$ mostra:
\begin{enumerate}
\item \textbf{Gradiente vertical:} bases maiores produzem mais atratores, consistente com $\Cyc(\varphi_{k,b-1})$ crescendo com $b-1$.
\item \textbf{Periodicidade colunar:} certos $k$ produzem sistematicamente menos ou mais atratores em todas as bases, refletindo a estrutura de $\gcd(k, b-1)$.
\item \textbf{Ressonâncias diagonais:} pontos quentes seguem trajetórias diagonais, sugerindo dependência de $k \bmod (b-1)$ ou $\gcd(k, \varphi(b-1))$.
\end{enumerate}

\section{Preditores Estruturais da Contagem de Atratores (Empíricos)}\label{sec:predictors}

O Teorema~\ref{thm:bdl} fixa os $p_i$, mas é silencioso sobre quantos atratores físicos compartilham $\gamma_i$. Relato preditores empíricos do lado da contagem (o excesso $\Delta$), enunciados como correlações, não formas fechadas, com sinalização onde se sobrepõem a estruturas conhecidas de números felizes.

\noindent\textbf{9.1 $\omega(b-1)$ é o preditor estrutural principal.} Com $\omega(m)$ o número de fatores primos distintos de $m = b-1$: Pearson $R = +0{,}436$ com $\Cyc$, $+0{,}419$ com $|C|$, $+0{,}391$ com $\Delta$. Para maximizar macroestados, tome $b-1$ com muitos primos distintos ($b-1 = 2\cdot3\cdot5 = 30 \Rightarrow b = 31$).

\medskip
\noindent\textbf{9.2 A estrutura $2$-ádica de $\gcd(k, b-1)$ controla o excesso em um único passo.} A correlação linear de $\gcd$ com $\Delta$ é negligenciável ($-0{,}05$), mascarando um padrão forte sob agrupamento:

\begin{center}
\begin{tabular}{@{}cc@{\qquad}cc@{}}
\toprule
$\gcd(k,b-1)$ & média $\Delta$ & $\gcd$ & média $\Delta$ \\
\midrule
$1$ & $14{,}46$ & $3$ & $14{,}15$ \\
$2$ & $8{,}79$ & $5$ & $16{,}55$ \\
$4$ & $6{,}90$ & $15$ & $26{,}88$ (máx.) \\
$8$ & $5{,}16$ & & \\
$16$ & $4{,}06$ & & \\
$32$ & $3{,}73$ (mín.) & & \\
\bottomrule
\end{tabular}
\end{center}

Potências de dois baixam $\Delta$; gcds ímpares (especialmente com vários primos ímpares) inflacionam-no.

\textbf{Cautela na leitura desta tabela.} Cada linha agrega pares de muitos módulos diferentes, e $v_2(\gcd(k,b-1))$ é grande só quando $b-1$ ele próprio carrega estrutura de potência de $2$, de modo que as linhas não são populações comparáveis. Dentro de um único módulo a contagem de ciclos não é nem monótona nem plana em $v_2(k)$: pelo Corolário~\ref{cor:notradical}, $m=41$ dá $\Cyc = 3, 4, 3, 6, 3$ para $k = 2, 4, 8, 16, 32$. A tendência decrescente na coluna das potências de dois é portanto um efeito de agregação entre módulos, e nenhuma lei por sistema deve ser lida nela. \texttt{scripts/analyze\_patterns.py} relata o desvio padrão intra-grupo e o número de módulos distintos em cada linha exatamente por esta razão. Um rascunho anterior foi além e afirmou um colapso geométrico graduado $\overline{\Delta}(v_2) \approx 12{,}68 \cdot (0{,}764)^{v_2}$; um posterior substituiu-o pela afirmação de que a contagem de ciclos é plana para todo $v_2 \ge 1$. Ambos estão retratados; ver Apêndice~B.1.

\textbf{Relação com resultados conhecidos.} Isto não é independente da literatura de números felizes: Grundman--Teeple \cite{GT07} mostram que existem sequências arbitrariamente longas de números $(e,b)$-felizes consecutivos precisamente quando $e-1$ não é divisível por $p-1$ para nenhum primo $p \mid b-1$, uma condição de coprimalidade do expoente contra a estrutura prima de $b-1$ do mesmo sabor do padrão acima. As seções 9.1--9.2 são um correspondente empírico e quantitativo para a família soma-de-dígitos-de-potências, não uma reivindicação independente desse círculo de ideias.

\medskip
\noindent\textbf{9.3 A paridade de $k$, não sua primalidade, amplifica atratores.} O contraste bruto replica (primo $k$: média $|C| = 31{,}82$ contra $18{,}26$ para composto, cerca de $74\%$ a mais) mas está confundido. Dois é o único primo par, enquanto cerca de metade dos compostos são pares, e $2 \mid k$ colapsa a parte $2$-ádica de todo grupo local de unidades (Proposição~\ref{prop:per}: $\kappa_k$ a retira), cortando tanto $\Cyc$ quanto $|C|$. Controlando a paridade numa grade exaustiva de $1{,}298$ pares ($2 \le k \le 60$, $3 \le b \le 24$):

\begin{center}
\begin{tabular}{@{}lccc@{}}
\toprule
classe de $k$ & $n$ & média $|C|$ & média $\Cyc$ \\
\midrule
primo (bruto) & $374$ & $17{,}27$ & $7{,}97$ \\
composto (bruto) & $924$ & $11{,}40$ & $4{,}82$ \\
primo ímpar & $352$ & $18{,}09$ & $8{,}25$ \\
composto ímpar & $286$ & $17{,}16$ & $7{,}35$ \\
ímpar (qualquer) & $638$ & $17{,}67$ & $7{,}85$ \\
par (qualquer) & $660$ & $8{,}66$ & $3{,}68$ \\
\bottomrule
\end{tabular}
\end{center}

Primos ímpares e compostos ímpares diferem de $0{,}93$ na média $|C|$, contra um desvio padrão intra-classe de cerca de $9{,}7$; ímpares e pares diferem de $9{,}01$. O efeito é a paridade. A correlação relacionada $R = -0{,}391$ entre a função número-de-divisores de $k$ e $\Cyc$ está sujeita ao mesmo confundidor, pois $k$ par tende a ter mais divisores. (\texttt{verification/audit/audit\_04\_parity\_confound.py}; Apêndice~B.4.)

\medskip
\noindent\textbf{9.4 A profundidade transiente escala com $k$.} A profundidade transiente máxima correlaciona $R = +0{,}540$ com $k$; a mais profunda observada: $133$ passos.

\section{Problemas em Aberto}

\begin{problem}[Fórmula exata]
Encontrar uma expressão em forma fechada ou um algoritmo eficiente para $|C(k,b)|$ além do Teorema~\ref{thm:lb}. Em particular, caracterizar o excesso de bifurcação $\Delta(k,b)$.
\end{problem}

\begin{problem}[Caracterização da justeza]\label{prob:tight}
Para quais pares $(k,b)$ vale a igualdade $|C(k,b)| = \Cyc(\varphi_{k,b-1})$? Os dados sugerem $b$ pequeno (especialmente $b = 2,3$) ou dinâmica modular trivial; a frequência de igualdade exata é ela própria sensível à faixa.
\end{problem}

\begin{problem}[Assíntotica do comprimento de órbita]
Seja $L(k,b,N)$ o comprimento máximo de órbita (transiente + ciclo) entre $n \le N$. Qual é sua taxa de crescimento? Vale $L(k,b,N) = O(\log N)$ para todo $(k,b)$ fixo?
\end{problem}

\begin{problem}[Estrutura do grafo de predecessores]
Para $(k,b)$ fixo e $A \in C(k,b)$, a árvore de bacia enraizada em $A$: qual é sua distribuição de graus? Segue uma lei de potência?
\end{problem}

\begin{problem}[Cota superior]
Provar uma cota superior complementar para $|C(k,b)|$; um candidato natural é o limiar de contração $N^*(k,b)$ do Lema~\ref{lem:contraction}.
\end{problem}

\begin{problem}[Teoria rigorosa da divisão]\label{prob:split}
Substituir a gaussiana numérica da \S\ref{sec:split} pelo termo de Fourier exato de Mellin--Perron para $S_b(n^k)$ e provar rigorosamente a oscilação quase-periódica da divisão por dígito; determinar se o enunciado já está implícito em \cite{delange,drmota-poly,drmota-grabner}.
\end{problem}

\begin{problem}[Lado da contagem do excesso]\label{prob:excess}
O lado modular da contagem está agora completamente fechado: tanto a contagem de pontos periódicos (Proposição~\ref{prop:per}) quanto a contagem de ciclos (Proposição~\ref{prop:cyc}) têm formas fechadas. Caracterizar o excesso físico $\Delta = |C| - \Cyc$, atratores extras da dinâmica de dígitos e não da estrutura modular. Empiricamente $\Delta$ responde à paridade de $k$ e a $\omega(b-1)$ (\S\ref{sec:predictors}), mas nenhuma forma fechada é conhecida.
\end{problem}

\begin{problem}[Complexidade da contagem de ciclos]
A Proposição~\ref{prop:cyc} avalia $\Cyc(\varphi_{k,m})$ a partir da fatoração de $m$ e de ordens multiplicativas módulo divisores de $\lambda(m)$, logo é polinomial dada a fatoração; a soma ingênua sobre $t$-uplas de ciclos locais não o é. Existe uma avaliação cujo custo é polinomial em $\log m$ e no número de fatores primos, sem enumerar $t$-uplas?
\end{problem}

\section{Conclusão}

Para a família $f_{k,b}(n) = S_b(n^k)$, a aritmética modular módulo $b-1$ determina, exatamente e elementarmente, tanto uma cota inferior para o número de atratores (Teorema~\ref{thm:lb}) quanto a massa total de bacia de cada assinatura de resíduos (Teorema~\ref{thm:bdl}, cujo conteúdo em janela finita é uma identidade exata entre inteiros, com a cota estreita $\min(|R_i|,b-2)/M$). Ambos foram verificados sem exceção ao longo da varredura exaustiva de $19{,}500$ sistemas. O lado modular da contagem está fechado em ambas as suas funções de contagem: a fórmula clássica de pontos periódicos da Proposição~\ref{prop:per}, e a fórmula de contagem de ciclos da Proposição~\ref{prop:cyc}, que é a que o Teorema~\ref{thm:lb} precisa, e que mostra que $\Cyc$, diferentemente de $\#\Per$, não é função de $\rad(k)$.

Esses resultados exatos quocientam a parte elementar da distribuição e isolam o objeto restante: a divisão intra-assinatura. Essa divisão não se estabiliza. Ao longo de comprimento de dígitos fixo, as massas de bacia individuais oscilam quase-periodicamente, em antifase entre atratores competidores, sempre somando o agregado exato $p_i$ (Observação~\ref{rem:result}). Isso vale para toda assinatura com excesso local positivo, inclusive as que um rascunho anterior classificou como convergentes. Um modelo de varredura gaussiana sem parâmetros reproduz a oscilação dentro do ruído de Monte Carlo. A oscilação é uma instância da clássica flutuação log-periódica de funções digitais (Delange; Drmota--Grabner). A contribuição do presente trabalho é uma lei agregada corretamente enunciada e computacionalmente verificada, uma forma fechada para a contagem de ciclos, e uma caracterização de por que a distribuição mais fina parece não existir para esta família.

\appendix

\section{Reprodutibilidade e Verificação Independente}

\textbf{Pipeline primário.} Tudo é reproduzível a partir do pacote \texttt{dspm} que acompanha este artigo. A varredura exaustiva por trás da \S8.2 é \texttt{scripts/sweep.py}, produzindo \texttt{data/sweeps/results\_k1-500\_b2-40\_*.jsonl.gz} (19.500 registros). O programa de divisão da \S\ref{sec:split} é \texttt{scripts/split\_scale.py} (medição) e \texttt{scripts/split\_predict.py} (o modelo sem parâmetros); o lado da contagem da \S\ref{sec:count} é \texttt{scripts/cycle\_structure.py}. Estatísticas da \S\ref{sec:predictors} são \texttt{scripts/analyze\_patterns.py}. Figuras: \texttt{paper/figures/split\_oscillation.svg}, \texttt{split\_predict\_overlay.svg}.

\begin{verbatim}
pip install -e ".[fast,dev]"
pytest                                        # suíte de testes no nível dos teoremas
python verification/verify_theorems.py        # rederivação independente
python scripts/cycle_structure.py             # ambas as formas fechadas vs. força bruta
python scripts/split_scale.py --k 3 --b 10 --d-max 90 --samples 60000
\end{verbatim}

\medskip
\noindent\textbf{Reverificação independente.} \texttt{verification/verify\_theorems.py} reimplementa a soma de dígitos, a região de aprisionamento, ambos os grafos funcionais e a estrutura modular do zero, sem importar nada de \texttt{dspm}, de modo que um bug no pacote não pode fazer um teorema parecer válido. Em $b \in [2,11]$, $k \in [1,15]$ relata:
\begin{itemize}
\item Lema~\ref{lem:inv} (invariância modular): $40{,}365$ testes, 0 falhas.
\item Teorema~\ref{thm:lb} ($|C| \ge \Cyc$): exaustivo sobre $150$ pares, 0 violações; igualdade exata em $19$ deles.
\item Lema~\ref{lem:sig} (a assinatura é um ciclo modular completo): 0 violações.
\item Proposição~\ref{prop:partition} (identidade inteira exata), sobre três janelas por par: 0 violações.
\item Teorema~\ref{thm:bdl}, contra tanto a cota enunciada quanto a estreita $\min(|R_i|,b-2)/M$: 0 violações; pior razão erro-cota $0{,}25$ para a primeira e $0{,}90$ para a segunda, que é a evidência de que a cota original não era estreita.
\item Proposição~\ref{prop:per}: forma fechada contra força bruta, $7{,}722$ pares, 0 discrepâncias.
\item $(k,b)=(2,10)$: atratores $\{1\},\{9\},\{13,16\}$ com proporções de bacia medidas $0{,}2223/0{,}3333/0{,}4443$ em $[1,3000]$, conferindo com $2/9, 3/9, 4/9$.
\item $(k,b)=(3,10)$: atratores $\{1\},\{8\},\{17\},\{18\},\{19,28\},\{26\},\{27\}$, conferindo com o Exemplo~\ref{ex:bif}.
\end{itemize}

A Proposição~\ref{prop:cyc} é checada à parte contra enumeração por força bruta para todos os $k \in [1,60]$, $m \in [1,259]$: $15{,}540$ pares, 0 discrepâncias.

\section{Errata}\label{app:errata}

Quatro afirmações do rascunho de julho de 2026 estavam erradas. Cada uma é registrada com o raciocínio que a produziu, porque em três dos quatro casos o erro estava no método de verificação e não na matemática, e esses modos de falha generalizam.

\medskip
\noindent\textbf{B.1. A contagem de ciclos modulares é plana para $v_2(k) \ge 1$.} \emph{Retirada.} A afirmação foi enunciada como consequência da Proposição~\ref{prop:per} e ``verificada para $m = 16, 17, 32, 37, 41, 64$''. Dois defeitos independentes. Primeiro, a Proposição~\ref{prop:per} conta pontos periódicos, e $\#\Per$ é de fato função de $\rad(k)$; a contagem de ciclos não o é, porque os comprimentos de ciclo são ordens multiplicativas $\ord_d(k)$ (Proposição~\ref{prop:cyc}, Corolário~\ref{cor:notradical}). Segundo, a computação de suporte agrupou $k$ pela parte ímpar e $v_2$ de $\gcd(k,\lambda(m))$ e comparou médias de grupo. Para os dois módulos não degenerados da lista os valores intra-grupo são $\{4,6,10\}$ para $m=37$ e $\{3,4,6\}$ para $m=41$; as médias concordam entre grupos de $v_2$ porque cada grupo contém o mesmo multiconjunto nas mesmas proporções. Os outros quatro módulos são degenerados: $\kappa_k$ colapsa a parte de unidades, $\#\Per = 2$, e a planaridade é vazia. A lição é que um teste comparando médias de grupo não detecta não-constância, e que quatro de seis casos de teste eram incapazes de falhar. Substituída pela Proposição~\ref{prop:cyc} e pelo Corolário~\ref{cor:notradical}.

\medskip
\noindent\textbf{B.2. As assinaturas $\{1\}$ e $\{8\}$ têm divisões convergentes.} \emph{Retirada.} O script de medição decidia convergência a partir da mudança na divisão ao longo de um passo de $D$, chamando-a estabilizada quando essa mudança caía abaixo do ruído de amostragem. Como as curvas são quase-periódicas com trechos longos e planos, a regra reporta convergência sempre que a janela amostrada se situa dentro de um deles, e a execução original parou em $D = 60$, dentro de tal trecho para dois dos três pontos fixos. Remedir em $4 \le D \le 90$ com $60{,}000$ amostras por faixa mostra os três pontos fixos ``desaparecendo'' retornando após longos silêncios, com picos de $0{,}027$ em $D = 79$ para $\{17\}$, $0{,}028$ em $D = 58$ para $\{8\}$, e $0{,}026$ em $D = 73$ para $\{1\}$; todas as sete curvas são não monótonas. Corrigida na \S7.4. A correção fortalece a afirmação empírica principal: toda densidade individual deixa de se assentar, não apenas a da assinatura $\{0\}$.

\medskip
\noindent\textbf{B.3. A cota em janela finita $(b-1)/M$ é ``saturada''.} \emph{Retirada.} Não é: ao longo da grade a pior razão entre erro real e cota é $0{,}25$. A prova do Teorema~\ref{thm:bdl} limita o desvio de uma classe de resíduos por $|R_i|$, mas os desvios ao longo de todas as $m$ classes somam zero, logo a constante estreita é $\min(|R_i|,\,m-|R_i|)$, e numa janela $M = qm + s$ é $\min(|R_i|, s) \le \min(|R_i|, b-2)$. A \S8.2 já exibia o hiato (pior erro $0{,}10$ contra uma cota de $0{,}40$) sem tirar a conclusão. Corrigida na \S1.3, \S8.2 e \S11; a cota estreita é verificada em \texttt{verification/audit/audit\_03\_bound\_sharpness.py}.

\medskip
\noindent\textbf{B.4. Expoentes primos amplificam atratores em $\approx 74\%$.} \emph{Reinterpretada.} O número replica, mas o contraste está confundido com a paridade de $k$, que é o motor real (Proposição~\ref{prop:per}). Comparação controlada na \S9.3.

\medskip
\noindent\textbf{Uma nota metodológica.} Três dos quatro erros sobreviveram porque a verificação foi executada contra uma estatística que não podia distinguir a afirmação de sua negação: uma comparação de médias onde a constância estava em questão, uma inclinação local onde um limite estava em questão, e uma desigualdade frouxa onde a estreiteza estava em questão. A contramedida adotada ao longo do código que acompanha o artigo é preferir identidades inteiras exatas onde elas existem, como na Proposição~\ref{prop:partition}, e, onde não existem, relatar amplitude, dispersão e piso de ruído ao lado de cada veredito.

\begin{thebibliography}{99}

\bibitem{guy} R.~K. Guy, ``Happy Numbers,'' in \emph{Unsolved Problems in Number Theory}, 3rd ed., Springer, 2004, pp.~358--359.

\bibitem{GT01} H.~G. Grundman and E.~A. Teeple, ``Generalized happy numbers,'' \emph{The Fibonacci Quarterly} \textbf{39} (2001), 462--466.

\bibitem{GT07} H.~G. Grundman and E.~A. Teeple, ``Sequences of consecutive happy numbers,'' \emph{Rocky Mountain J. Math.}; and ``Sequences of generalized happy numbers with small bases,'' \emph{J. Integer Sequences} \textbf{10} (2007), Article 07.1.8.

\bibitem{hardy-wright} G.~H. Hardy and E.~M. Wright, \emph{An Introduction to the Theory of Numbers}, 6th ed., Oxford University Press, 2008.

\bibitem{survey} E.~Hasson et al.\ (eds.), ``Happy Numbers, Happy Functions, and Their Variations: A Survey,'' \emph{La Matematica} (Springer), 2021.

\bibitem{gilmer} J.~Gilmer, ``On the density of happy numbers,'' \emph{Integers} \textbf{13} (2013), \#A48 (arXiv:1110.3836).

\bibitem{mauduit-rivat} C.~Mauduit and J.~Rivat, ``La somme des chiffres des carr\'es,'' \emph{Acta Mathematica} \textbf{203} (2009), 107--148.

\bibitem{dmr} M.~Drmota, C.~Mauduit, and J.~Rivat, ``The sum of digits of polynomial values in arithmetic progressions,'' and related local-limit results for $S_b(n^k)$.

\bibitem{silverman} J.~H. Silverman, \emph{The Arithmetic of Dynamical Systems}, GTM 241, Springer, 2007.

\bibitem{somer-krizek} L.~Somer and M.~K\v{r}\'i\v{z}ek, ``On a connection of number theory with graph theory,'' \emph{Czechoslovak Math. J.} \textbf{54} (2004), 465--485; and ``Structure of digraphs associated with quadratic congruences with composite moduli,'' \emph{Discrete Math.} \textbf{306} (2006).

\bibitem{chou-shparlinski} W.-S. Chou and I.~E. Shparlinski, ``On the cycle structure of repeated exponentiation modulo a prime,'' \emph{J. Number Theory} \textbf{107} (2004), 345--356.

\bibitem{delange} H.~Delange, ``Sur la fonction sommatoire de la fonction `somme des chiffres','' \emph{Enseign. Math.} \textbf{21} (1975), 31--47.

\bibitem{drmota-poly} M.~Drmota, ``The sum of digits function of polynomial sequences,'' TU Wien. Gaussian distribution and periodic fluctuations of $S_b(P(n))$, including $P(n) = n^k$.

\bibitem{drmota-grabner} M.~Drmota and P.~Grabner, \emph{Analysis of Digital Functions and Applications}, Encyclopedia of Mathematics and Its Applications, Cambridge University Press, 2010.

\bibitem{devaney} R.~L. Devaney, \emph{An Introduction to Chaotic Dynamical Systems}, 2nd ed., Westview Press, 2003.

\bibitem{porges} A.~Porges, ``A set of eight numbers,'' \emph{American Mathematical Monthly} \textbf{52} (1945), 379--382.

\end{thebibliography}

\end{document}
